\chapter{Justificativa}
\label{justificativa}
A transição para um modelo computacional especializado não visa apenas a digitalização de documentos, mas a garantia da integridade e da agilidade institucional por meio da Engenharia de Software\cite{}.

\section{Dimensões Complementares}
A justificativa para o desenvolvimento deste sistema estrutura-se em duas frentes complementares. No âmbito administrativo, busca-se a otimização de recursos e a redução de retrabalho. No âmbito acadêmico, o foco reside na mitigação de riscos e na redução da carga burocrática sobre o estudante. Juntas, essas dimensões fundamentam a aplicação de soluções automatizadas como resposta a processos manuais ineficientes. 

\subsection{Âmbito Administrativo}

Do ponto de vista administrativo, a coordenação do curso opera com recursos humanos limitados e precisa analisar individualmente cada barema submetido. O volume de documentos incorretos que retornam para correção representa um custo operacional real: cada revisão exige nova comunicação com o discente, novo prazo de reenvio e nova análise pelo coordenador. 

Nesse cenário, a automação do fluxo não é apenas uma conveniência, mas uma aplicação direta de princípios consolidados na Engenharia de Software. A automação minimiza a intervenção humana em tarefas repetitivas, reduzindo o risco de erros operacionais e aumentando a eficiência \cite{abidemi2024role}. Assim, a indexação de páginas e o cálculo de cargas horárias são tarefas determinísticas e com regras bem definidas, beneficiando-se diretamente dessa abordagem para otimizar a capacidade de atendimento do colegiado.

\subsection{Meio Acadêmico}

Do ponto de vista acadêmico, o processo atual impõe ao discente um ônus desproporcional à natureza da tarefa. Montar um barema não é uma atividade de aprendizado — é uma tarefa burocrática de baixa complexidade cognitiva, mas de alto potencial de erro quando executada manualmente. Exigir que o aluno dedique tempo e atenção a contar páginas e somar horas em planilhas é desviar energia que poderia ser investida nas próprias atividades complementares que o processo visa reconhecer. Além disso, erros nessa etapa têm consequências concretas: atrasos no aproveitamento das AC podem comprometer o cumprimento dos prazos para a colação de grau, impactando diretamente o planejamento acadêmico e profissional do estudante.

\section{Justificativa Técnica}

Por fim, a justificativa técnica do projeto reside no fato de que a solução proposta não é apenas um agregador de arquivos, mas um sistema que codifica as regras de negócio do COLCIC. Isso significa que as validações de carga horária máxima por categoria, a geração automática da tabela de indexação e a unificação dos comprovantes ocorrem de forma integrada e transparente para o usuário. O resultado é um documento gerado com garantia de conformidade regulatória — algo que nenhuma ferramenta genérica hoje disponível é capaz de oferecer para esse contexto específico.
